\documentclass[10pt, letterpaper]{article}

% Packages:
\usepackage[
	ignoreheadfoot, % set margins without considering header and footer
	top=1 cm, % seperation between body and page edge from the top
	bottom=1 cm, % seperation between body and page edge from the bottom
	left=1 cm, % seperation between body and page edge from the left
	right=1 cm, % seperation between body and page edge from the right
	footskip=1 cm, % seperation between body and footer
	% showframe % for debugging
]{geometry} % for adjusting page geometry
\usepackage{titlesec} % for customizing section titles
\usepackage[dvipsnames]{xcolor} % for coloring text
\definecolor{primaryColor}{RGB}{0, 0, 0} % define primary color
\usepackage{enumitem} % for customizing lists
\usepackage{etoolbox} % for conditional statements
\usepackage{amsmath} % for math
\usepackage[
	pdftitle={Curriculo de Marllon Rodrigues},
	pdfauthor={Marllon Rodrigues},
	pdfcreator={Marllon Rodrigues},
	colorlinks=true,
	urlcolor=primaryColor
]{hyperref} % for links, metadata and bookmarks
\usepackage[pscoord]{eso-pic} % for floating text on the page
\usepackage{calc} % for calculating lengths
\usepackage{bookmark} % for bookmarks
\usepackage{lastpage} % for getting the total number of pages
\usepackage{changepage} % for one column entries (adjustwidth environment)
\usepackage{paracol} % for two and three column entries
\usepackage{needspace} % for avoiding page brake right after the section title
\usepackage{iftex} % check if engine is pdflatex, xetex or luatex

% Ensure that generate pdf is machine readable/ATS parsable:
\ifPDFTeX
\input{glyphtounicode}
\pdfgentounicode=1
\usepackage[T1]{fontenc}
\usepackage[utf8]{inputenc}
\usepackage{lmodern}
\fi

\usepackage{charter}

% Some settings:
\raggedright
\AtBeginEnvironment{adjustwidth}{\partopsep0pt} % remove space before adjustwidth environment
\pagestyle{empty} % no header or footer
\setcounter{secnumdepth}{0} % no section numbering
\setlength{\parindent}{0pt} % no indentation
\setlength{\topskip}{0pt} % no top skip
\setlength{\columnsep}{0.15cm} % set column seperation
\pagenumbering{gobble} % no page numbering

\titleformat{\section}{\needspace{4\baselineskip}\bfseries\large}{}{0pt}{}[
\vspace{1pt}
\titlerule]

\titlespacing{\section}{
% left space:
-1pt }{
% top space:
0.3 cm }{
% bottom space:
0.2 cm } % section title spacing

\renewcommand{\labelitemi}{$\vcenter{\hbox{\small$\bullet$}}$} % custom bullet points
\newenvironment{highlights}{ \begin{itemize}[ topsep=0.10 cm, parsep=0.10 cm, partopsep=0pt,
itemsep=0pt, leftmargin=0 cm + 10pt ] }{ \end{itemize} } % new environment for highlights

\newenvironment{highlightsforbulletentries}{ \begin{itemize}[ topsep=0.10 cm,
parsep=0.10 cm, partopsep=0pt, itemsep=0pt, leftmargin=10pt ] }{ \end{itemize} } % new environment for highlights for bullet entries

\newenvironment{onecolentry}{ \begin{adjustwidth}{ 0 cm + 0.00001 cm }{ 0 cm + 0.00001 cm }
}{ \end{adjustwidth} } % new environment for one column entries

\newenvironment{twocolentry}[2][]{ \onecolentry \def\secondColumn{#2} \setcolumnwidth{\fill, 4.5 cm}
\begin{paracol}{2} }{ \switchcolumn \raggedleft \secondColumn \end{paracol}
\endonecolentry } % new environment for two column entries

\newenvironment{threecolentry}[3][]{ \onecolentry 
\def\thirdColumn{#3} \setcolumnwidth{, \fill, 4.5 cm}
\begin{paracol}{3} {\raggedright #2} \switchcolumn }{ \switchcolumn \raggedleft \thirdColumn
\end{paracol} \endonecolentry } % new environment for three column entries

\newenvironment{header}{
\setlength{\topsep}{0pt}
\par\kern\topsep
\centering
\linespread{1.5} }{ \par\kern\topsep } % new environment for the header

\newcommand{\placelastupdatedtext}{% \placetextbox{<horizontal pos>}{<vertical pos>}{<stuff>}
\AddToShipoutPictureFG*{% Add <stuff> to current page foreground
\put( \LenToUnit{\paperwidth-2 cm-0 cm+0.05cm}, \LenToUnit{\paperheight-1.0 cm} ){\vtop{{\null}\makebox[0pt][c]{ \small\color{gray}\textit{Última atualização em Novembro de 2025}\hspace{\widthof{Última atualização em Novembro de 2025}} }}}%
}%
}%

% save the original href command in a new command:
\let\hrefWithoutArrow\href

% new command for external links:

\begin{document}
	\newcommand{\AND}{\unskip \cleaders\copy\ANDbox\hskip\wd\ANDbox \ignorespaces }
	\newsavebox{\ANDbox}
	\sbox{\ANDbox}{$|$}

	\begin{header}
		\fontsize{25 pt}{25 pt}\selectfont Marllon Rodrigues

        \fontsize{18 pt}{18 pt}\selectfont Engenheiro de Software \& Automação

		\vspace{5 pt}

		\normalsize
		\mbox{\hrefWithoutArrow{mailto:marllonrsantos@gmail.com}{marllonrsantos@gmail.com}}%
		\kern 5.0 pt%
		\AND%
		\kern 5.0 pt%
		\mbox{\hrefWithoutArrow{https://linkedin.com/in/rodriguesmarllon}{linkedin.com/in/rodriguesmarllon}}%
		\kern 5.0 pt%
        \AND%
		\kern 5.0 pt%
		\mbox{\hrefWithoutArrow{https://github.com/rodriguesmarllon}{github.com/rodriguesmarllon}}%
	\end{header}

	\section{Competências Técnicas}

    \textbf{Desenvolvimento de Software:} Java, C\# (.NET), Python (FastAPI), Node.js (NestJS), TypeScript, Vue.js, React, Next.js, Prisma, Supabase, PostgreSQL, Docker, CI/CD (GitHub Actions), Git.
    
    \textbf{Automação Industrial:} PLC (Siemens S7-1200/1500, Rockwell Control/Compact/MicroLogix, Schneider), HMI/SCADA (AVEVA System Platform, InTouch, Historian, FactoryTalk), SQL Server, Redes Industriais, Integração IT/OT.

	\section{Experiência}
	
	\begin{twocolentry}
		{ Set 2025 - Atual } \textbf{Co-Founder \& Tech Lead}, Startup de Tecnologia (Em Desenvolvimento)
	\end{twocolentry}
	
	\vspace{0.10 cm}
	\begin{onecolentry}
		\begin{highlights}
			\item \textbf{Liderança de Produto:} Arquitetura de dois produtos SaaS B2B/B2C: um ERP para Gestão Municipal (encaminhamentos, NF, agendas) e uma Plataforma Financeira 360º (Investimentos, Cripto e Sucessão Patrimonial).
            \item \textbf{DevOps \& Qualidade:} Configuração de pipeline rigoroso de CI/CD (GitHub Actions) para deploy automatizado (Vercel/Render) com execução de testes (Jest) em cada Commit e Pull Request.
            \item \textbf{Mentoria \& Capacitação:} Capacitação de desenvolvedor júnior, que evoluiu para uma atuação full stack autônoma na resolução de Issues e entrega de features após mentoria técnica intensiva.
			\item \textbf{Stack Tecnológica:} Utilização de Next.js, NestJS, Prisma e Supabase para criar soluções escaláveis e de alta performance.
		\end{highlights}
	\end{onecolentry}

    \vspace{0.2 cm}

    \begin{twocolentry}
		{ Jun 2025 - Atual } \textbf{Desenvolvedor Full Stack}, Bunkerchain (Singapura/Remoto)
	\end{twocolentry}
	
	\vspace{0.10 cm}
	\begin{onecolentry}
		\begin{highlights}
            \item \textbf{Desenvolvimento de Produto (Touch and Sail):} Atuação no desenvolvimento do sistema "Touch and Sail" (e-BDN), uma solução de \textit{Digital Bunkering} para gestão de operações marítimas críticas em tempo real.
			\item \textbf{Frontend Complexo (Vue.js):} Desenvolvimento diário de interfaces dinâmicas e formulários para inserção de dados operacionais, garantindo alta responsividade e visualização de dados em tempo real.
            \item \textbf{Backend \& Integração:} Consumo intensivo de APIs Java e realização de manutenção corretiva (bug fix) no backend para assegurar a integridade dos dados e a comunicação fluida entre front e back.
            \item \textbf{Processos Internacionais:} Participação ativa em rituais ágeis (Dailies, Demos) em inglês, suporte técnico direto a clientes internacionais e gestão de bugs/tickets via Mantis.
		\end{highlights}
	\end{onecolentry}
	
	\vspace{0.2 cm}
	
	\begin{twocolentry}
		{ Ago 2024 - Mai 2025 } \textbf{Desenvolvedor de Sistemas}, MAIMA Tecnologia
	\end{twocolentry}
	
	\vspace{0.10 cm}
	\begin{onecolentry}
		\begin{highlights}
            \item \textbf{Automação SAP (RPA):} Desenvolvimento de robôs em Python para automação de 4-6 processos críticos diários no SAP, reduzindo o tempo de execução manual de 8 horas para apenas 2 minutos.
			\item \textbf{Manutenção Preditiva (Forecast):} Implementação de sistema em C\# para previsão de paradas de máquina com base na análise da qualidade do produto (cerdas dentais), antecipando falhas críticas.
            \item \textbf{Arquitetura de Microsserviços:} Segregação de responsabilidades em microsserviços distribuídos (Análise de Imagem, Dados, BFF), garantindo escalabilidade e manutenção simplificada.
            \item \textbf{Suporte \& Chatbot:} Contribuição no desenvolvimento de Chatbot com Adaptive Cards para suporte operacional ao sistema de qualidade.
		\end{highlights}
	\end{onecolentry}
	
	\vspace{0.2 cm}
	
	\begin{twocolentry}
		{ Jun 2022 - Ago 2024 } \textbf{Técnico em Automação}, S.Build Automação
	\end{twocolentry}
	
	\vspace{0.10 cm}
	\begin{onecolentry}
		\begin{highlights}
			\item \textbf{Automação de Dados (Python/MES):} Criação de scripts em Python para gerar Queries SQL dinâmicas, atualizando em massa milhares de registros no AVEVA MES, tarefa inviável manualmente.
            \item \textbf{Engenharia de Campo \& Startups:} Atuação intensiva em viagens para comissionamento de obras \textit{in loco}. Especialista na implementação de ecossistemas \textbf{AVEVA System Platform}, desenvolvimento de interfaces \textbf{InTouch} e configuração de \textbf{Historian} em ambientes críticos.
            \item \textbf{Tech Stack Industrial:} Projetos entregues com PLCs Siemens (S7-1200/1500), Rockwell (Control/Compact/MicroLogix) e Schneider. Configuração de redes (Profinet, Ethernet/IP, Modbus, OPC) e fluxos com \textbf{AVEVA WorkTasks}.
            \item \textbf{Infraestrutura \& Retrofit:} Migração de legado (1995 $\rightarrow$ 2023) e retrofit de máquinas sem documentação (engenharia reversa). Upgrade de 20+ servidores e gestão via ThinManager.
		\end{highlights}
	\end{onecolentry}

    \vspace{0.2 cm}

    \begin{twocolentry}
		{ Set 2021 - Fev 2023 } \textbf{Competidor WorldSkills (Mecatrônica)}, SENAI Roberto Simonsen
	\end{twocolentry}
	
	\vspace{0.10 cm}
	\begin{onecolentry}
		\begin{highlights}
            \item \textbf{Alta Performance \& Conquistas:} Selecionado como \textbf{Melhor Aluno (Top 1)} do SENAI Roberto Simonsen para compor a equipe de elite. A equipe sagrou-se \textbf{Campeã Estadual} (São Paulo Skills).
            \item \textbf{Treinamento de Elite:} Rotina de 8-10 horas diárias focada em excelência operacional, montagem mecânica/elétrica de precisão e programação sob pressão de tempo.
			\item \textbf{Tecnologia Festo MPS:} Montagem, comissionamento e programação de estações industriais complexas (\textit{Distributing, Sorting, Testing, Handling/Robot}), simulando processos reais de separação, medição e envase (Filling).
            \item \textbf{Indústria 4.0:} Integração avançada de sensores inteligentes via \textbf{IO-Link} e redes industriais, garantindo comunicação eficiente entre sensores e PLCs Siemens S7-1500.
		\end{highlights}
	\end{onecolentry}
	
    \section{Certificações}
    
    \begin{highlights}
        \item \textbf{CSI AVEVA} Manufacturing Execution System (MES) 2020 - Operations \& Performance Exams
        \item \textbf{CSI AVEVA} Application Server 2023 Exam
        \item \textbf{Formação .NET Developer} (DIO - 93h)
        \item \textbf{Java COMPLETO Programação Orientada a Objetos + Projetos} (Udemy - 54.5h)
        \item \textbf{Formação Node.js Fundamentals} (DIO - 30h)
    \end{highlights}

	\section{Educação}
	
	\begin{twocolentry}
		{2023 - 2025} \textbf{UNINTER}, Análise e Desenvolvimento de Sistemas
	\end{twocolentry}
	
	\begin{twocolentry}
		{2021 -  2022} \textbf{SENAI Roberto Simonsen}, Técnico em Redes de Computadores
	\end{twocolentry}

    \begin{twocolentry}
		{2020 -  2022} \textbf{SENAI Roberto Simonsen}, Eletricista de Manutenção Eletroeletrônica
	\end{twocolentry}
	
	\section{Idiomas}
	
	\begin{onecolentry}
		\begin{highlights}
			\item Inglês (C1 - Avançado)
			\item Português (Nativo)
		\end{highlights}
	\end{onecolentry}
	
\end{document}